\section{Do the perturbations transfer across models?}
\label{sec:transfer}

A single-model ASR answers a narrow question: how much does \emph{this} model cost
an attacker who can query it. That is the right question for a targeted campaign
and the wrong one for a defender choosing a model, because an attacker who cannot
query the deployed model will instead build perturbations against something they
can download. If those perturbations do not survive the move, the weakness
measured in \S\ref{sec:results} belongs to three specific checkpoints rather than
to the task.

\paragraph{Protocol.} We generate no new perturbations. Every text produced by
the three attacks --- character, word and saliency --- is replayed, unchanged,
against every other model, with the same prompt, the same answer parser and
temperature zero. The saliency rows are the ones worth having: that edit was
chosen against the source model's own measured probabilities, so its cell asks
whether the choice survives a different decision boundary, where a random edit
only asks whether the model is noisy in the same way as its neighbour.
Two counting rules make the resulting numbers comparable with the single-model
ones:

\begin{itemize}
\item \textbf{The target owns the denominator.} A perturbation ``transfers'' only
if the target model read the \emph{original} text as malicious and the
\emph{edited} text as benign. A message the target never flagged cannot be
evaded, and scoring it as a failed transfer would punish the target for being
wrong on the clean text --- the opposite of robustness.
\item \textbf{Two subsets, one tabulated.} The \emph{applied} subset contains
every perturbation the source actually applied; its denominator is how often the
target caught those originals, so it measures how well the perturbation
\emph{family} travels. The \emph{flipped-source} subset keeps only the
perturbations that fooled the source --- the classic transfer matrix --- and it is
what the literature reports, but at this scale it holds one to six candidates per
cell, so its cells carry no denominator worth an interval. We tabulate the applied
subset and keep the classic one in the artefact.
\end{itemize}

\paragraph{The faithfulness check.} Replaying a model's own perturbations must
reproduce its single-model ASR exactly --- same text, same prompt, deterministic
decoder --- and a mismatch there would invalidate every off-diagonal cell. We run
that replay as part of the same pass (\texttt{run\_transfer.py --include-self}),
and we compare it against the attack tables rather than against itself: the
stored flip flag from the attack run is the claim, the replayed verdict is the
evidence. Every diagonal cell agrees. This is what makes the off-diagonal cells a
measurement: the replay path is the same code, prompt and parser that produced
Table~\ref{tab:asr-word}, so a transferred flip cannot come from the test setup.
It is a weaker statement than it looks, and worth saying plainly: agreement shows
the replay is faithful, not that the decoder is deterministic in general ---
greedy decoding is reproducible on one GPU build, and a cell could drift on
another.

\IfFileExists{tables/t5_transfer.tex}{% Generated by src/make_paper_tex.py from results/transfer_matrix.csv.
On the applied subset 22 of 36 off-diagonal cells have a nonzero transfer rate, and every cell has a denominator worth an interval.
Among the cells with a denominator worth an interval, the highest rate is SmolLM2-1.7B-Instruct $\rightarrow$ Phi-3.5-mini-instruct on Phishing email under the saliency attack: 5 of 120 caught originals flipped (4.2\%, 95\% CI 1.8--9.4).
The self-replay is the check that makes the rest of the table readable: each model's own perturbations went through the same path a second time, and its flip count is compared with the one the attack recorded. 18 diagonal cells covering 3{,}080 rows were replayed, and every one reproduces the single-model ASR exactly.
}{}

Table~\ref{tab:transfer} reports the applied subset only, for all three attack
modes; the classic flipped-source matrix lives in the artefact, where its small
denominators cannot be mistaken for a rate.

\begin{table*}[t]
\centering
\small
\caption{Cross-model transfer on the applied subset: every perturbation the source produced is replayed on the target, and the denominator is the TARGET's own caught originals, because a perturbation cannot transfer on a message the target never flagged. No diagonal is tabulated: replaying a model's own perturbations reproduces its attack run and would only reprint Table~\ref{tab:asr-word}. That replay was run as a faithfulness check and its diagonal appears in the prose below, not as a result. Cells with fewer than ten caught originals would be marked $\dagger$; none in this table needed it.}
\label{tab:transfer}
\begin{tabular}{llllcccc}
\toprule
\textbf{Mode} & \textbf{Task} & \textbf{Source} & \textbf{Target} & \textbf{Caught} & \textbf{Flips} & \textbf{ASR} & \textbf{95\% CI} \\
\midrule
 char & SMS smishing & Phi-3.5-mini-instruct & Qwen2.5-3B-Instruct & 163 & 0 & 0.0\% & 0.0--2.3 \\
  &  & Phi-3.5-mini-instruct & SmolLM2-1.7B-Instruct & 160 & 0 & 0.0\% & 0.0--2.3 \\
  &  & Qwen2.5-3B-Instruct & Phi-3.5-mini-instruct & 163 & 1 & 0.6\% & 0.1--3.4 \\
  &  & Qwen2.5-3B-Instruct & SmolLM2-1.7B-Instruct & 179 & 0 & 0.0\% & 0.0--2.1 \\
  &  & SmolLM2-1.7B-Instruct & Phi-3.5-mini-instruct & 160 & 1 & 0.6\% & 0.1--3.5 \\
  &  & SmolLM2-1.7B-Instruct & Qwen2.5-3B-Instruct & 179 & 0 & 0.0\% & 0.0--2.1 \\
 char & Phishing email & Phi-3.5-mini-instruct & Qwen2.5-3B-Instruct & 119 & 0 & 0.0\% & 0.0--3.1 \\
  &  & Phi-3.5-mini-instruct & SmolLM2-1.7B-Instruct & 120 & 0 & 0.0\% & 0.0--3.1 \\
  &  & Qwen2.5-3B-Instruct & Phi-3.5-mini-instruct & 119 & 4 & 3.4\% & 1.3--8.3 \\
  &  & Qwen2.5-3B-Instruct & SmolLM2-1.7B-Instruct & 189 & 0 & 0.0\% & 0.0--2.0 \\
  &  & SmolLM2-1.7B-Instruct & Phi-3.5-mini-instruct & 120 & 2 & 1.7\% & 0.5--5.9 \\
  &  & SmolLM2-1.7B-Instruct & Qwen2.5-3B-Instruct & 189 & 1 & 0.5\% & 0.1--2.9 \\
\midrule
 word & SMS smishing & Phi-3.5-mini-instruct & Qwen2.5-3B-Instruct & 161 & 1 & 0.6\% & 0.1--3.4 \\
  &  & Phi-3.5-mini-instruct & SmolLM2-1.7B-Instruct & 158 & 1 & 0.6\% & 0.1--3.5 \\
  &  & Qwen2.5-3B-Instruct & Phi-3.5-mini-instruct & 162 & 3 & 1.8\% & 0.6--5.3 \\
  &  & Qwen2.5-3B-Instruct & SmolLM2-1.7B-Instruct & 178 & 0 & 0.0\% & 0.0--2.1 \\
  &  & SmolLM2-1.7B-Instruct & Phi-3.5-mini-instruct & 160 & 1 & 0.6\% & 0.1--3.5 \\
  &  & SmolLM2-1.7B-Instruct & Qwen2.5-3B-Instruct & 178 & 4 & 2.2\% & 0.9--5.6 \\
 word & Phishing email & Phi-3.5-mini-instruct & Qwen2.5-3B-Instruct & 117 & 0 & 0.0\% & 0.0--3.2 \\
  &  & Phi-3.5-mini-instruct & SmolLM2-1.7B-Instruct & 118 & 0 & 0.0\% & 0.0--3.1 \\
  &  & Qwen2.5-3B-Instruct & Phi-3.5-mini-instruct & 119 & 4 & 3.4\% & 1.3--8.3 \\
  &  & Qwen2.5-3B-Instruct & SmolLM2-1.7B-Instruct & 188 & 0 & 0.0\% & 0.0--2.0 \\
  &  & SmolLM2-1.7B-Instruct & Phi-3.5-mini-instruct & 120 & 3 & 2.5\% & 0.9--7.1 \\
  &  & SmolLM2-1.7B-Instruct & Qwen2.5-3B-Instruct & 188 & 1 & 0.5\% & 0.1--2.9 \\
\midrule
 saliency & SMS smishing & Phi-3.5-mini-instruct & Qwen2.5-3B-Instruct & 163 & 1 & 0.6\% & 0.1--3.4 \\
  &  & Phi-3.5-mini-instruct & SmolLM2-1.7B-Instruct & 160 & 0 & 0.0\% & 0.0--2.3 \\
  &  & Qwen2.5-3B-Instruct & Phi-3.5-mini-instruct & 154 & 2 & 1.3\% & 0.4--4.6 \\
  &  & Qwen2.5-3B-Instruct & SmolLM2-1.7B-Instruct & 163 & 1 & 0.6\% & 0.1--3.4 \\
  &  & SmolLM2-1.7B-Instruct & Phi-3.5-mini-instruct & 160 & 0 & 0.0\% & 0.0--2.3 \\
  &  & SmolLM2-1.7B-Instruct & Qwen2.5-3B-Instruct & 179 & 4 & 2.2\% & 0.9--5.6 \\
 saliency & Phishing email & Phi-3.5-mini-instruct & Qwen2.5-3B-Instruct & 119 & 1 & 0.8\% & 0.1--4.6 \\
  &  & Phi-3.5-mini-instruct & SmolLM2-1.7B-Instruct & 120 & 0 & 0.0\% & 0.0--3.1 \\
  &  & Qwen2.5-3B-Instruct & Phi-3.5-mini-instruct & 113 & 4 & 3.5\% & 1.4--8.8 \\
  &  & Qwen2.5-3B-Instruct & SmolLM2-1.7B-Instruct & 167 & 1 & 0.6\% & 0.1--3.3 \\
  &  & SmolLM2-1.7B-Instruct & Phi-3.5-mini-instruct & 120 & 5 & 4.2\% & 1.8--9.4 \\
  &  & SmolLM2-1.7B-Instruct & Qwen2.5-3B-Instruct & 189 & 5 & 2.6\% & 1.1--6.0 \\
\bottomrule
\end{tabular}
\end{table*}

